\documentclass[a4paper,11pt]{article}
\usepackage[T1]{fontenc}
\usepackage[utf8]{inputenc}
\usepackage[french]{babel}
\usepackage{geometry}
\geometry{top=1.5cm, bottom=1.8cm, left=1.8cm, right=1.8cm}
\usepackage{xcolor}
\renewcommand{\familydefault}{\sfdefault}
\definecolor{Couleur1}{HTML}{8b1f30}
\definecolor{Couleur6}{HTML}{CF6876}
\definecolor{Couleur7}{HTML}{A5C03F}
\definecolor{Couleur8}{HTML}{D36740}
\usepackage{amsmath, amssymb}
\usepackage{array, colortbl}
\usepackage{tikz}
\usepackage{fontawesome5}
\usepackage{mdframed}
\usepackage{fancyhdr}
\usepackage{enumitem}
\usepackage{multicol}
\usepackage{rotating}
\setlength{\parindent}{0pt}
\newcommand{\bl}[1][4cm]{\,\makebox[#1]{\dotfill}\,}
\newcommand{\lignepoints}{\noindent\dotfill\vspace{0.3em}\par}
\newcommand{\coeur}{\textcolor{Couleur8}{\faHeart}\ }
\newcommand{\coeurrappel}{\textcolor{Couleur6}{\faHeart}\ }
\pagestyle{fancy}
\fancyhf{}
\renewcommand{\headrulewidth}{0pt}
\fancyfoot[L]{\small Mme Bertail}
\fancyfoot[C]{\small \textcolor{Couleur6}{Interro de le\c{c}on n\degre~4 -- Ch 1 a 4}}
\fancyfoot[R]{\small \thepage}
\begin{document}
\begin{mdframed}[linecolor=Couleur1,linewidth=2pt,topline=true,bottomline=true,leftline=false,rightline=false,innertopmargin=8pt,innerbottommargin=8pt,backgroundcolor=Couleur1!5]
\begin{center}
{\Large\bfseries\textcolor{Couleur1}{Interrogation de le\c{c}on n\degre~4}}\\[4pt]
{\large\textcolor{Couleur8}{Chapitre 4 --- Angles (1)}}
\end{center}
\end{mdframed}
\vspace{0.4cm}
\noindent\textbf{Nom :} \dotfill\quad\textbf{Pr\'enom :} \dotfill\quad\textbf{Classe :} \dotfill\quad\textbf{Date :} \dotfill
\vspace{0.5cm}

\vspace{0.3cm}
\noindent{\textcolor{Couleur8}{\rule{\textwidth}{0.5pt}}}\\
{\small\textit{\textcolor{Couleur8}{Nouveau --- Chapitre 4}}}
\vspace{0.1cm}

\noindent \coeur \textbf{Compl\'eter les d\'efinitions.}
\begin{mdframed}[linecolor=Couleur6,linewidth=4pt,topline=false,bottomline=false,rightline=false,backgroundcolor=Couleur6!10,innerleftmargin=8pt,innertopmargin=6pt,innerbottommargin=6pt]
Un \textbf{angle aigu} est un angle \bl[6cm]\\[4pt]
Un \textbf{angle obtus} est un angle \bl[6cm]
\end{mdframed}

\noindent \coeur \textbf{Compl\'eter les d\'efinitions.} Deux angles sont dits :
\begin{mdframed}[linecolor=Couleur6,linewidth=4pt,topline=false,bottomline=false,rightline=false,backgroundcolor=Couleur6!10,innerleftmargin=8pt,innertopmargin=6pt,innerbottommargin=6pt]
\begin{itemize}[label=\textcolor{Couleur6}{$\bullet$},nosep,itemsep=4pt]
\item \textbf{suppl\'ementaires} lorsqu'ils \dotfill
\item \textbf{adjacents} lorsqu'ils \dotfill
\item \textbf{oppos\'es par le sommet} lorsqu'ils \dotfill\dotfill
\end{itemize}
\end{mdframed}

\noindent \coeur \textbf{Compl\'eter les \'etapes} pour mesurer un angle avec un rapporteur.
\begin{mdframed}[linecolor=Couleur1,linewidth=4pt,topline=false,bottomline=false,rightline=false,backgroundcolor=Couleur1!5,innerleftmargin=8pt,innertopmargin=6pt,innerbottommargin=6pt]
\begin{enumerate}[label=\textcolor{Couleur1}{\textbf{\arabic*.}},nosep,itemsep=4pt]
\item Placer le \textbf{centre} du rapporteur sur le \bl[3cm] de l'angle.
\item Aligner un \textbf{c\^ot\'e} de l'angle avec un \bl[1.5cm] du rapporteur.
\item Rep\'erer o\`u \bl[4cm] coupe le demi-cercle du rapporteur.
\item Lire la mesure en choisissant \bl[5cm].
\end{enumerate}
\end{mdframed}

\vspace{0.3cm}
\noindent{\textcolor{Couleur8}{\rule{\textwidth}{0.5pt}}}\\
{\small\textit{\textcolor{Couleur8}{Rappel --- Chapitre 3}}}
\vspace{0.1cm}

\noindent \coeurrappel $<$ se lit \bl[2.5cm] ; $>$ se lit \bl[2.5cm] ; $\leq$ se lit \bl[3cm] ; $\geq$ se lit \bl[3cm].
\end{document}